% !TeX root = ../../main.tex

\section{Big Data Analytics}

\subsection{¿Qué es Big Data Analytics?}

Big Data Analytics es el proceso de recopilar, preparar, almacenar y analizar
conjuntos de datos demasiado grandes, rápidos o variados para ser tratados
eficientemente con herramientas tradicionales.

Puede trabajar con datos estructurados, como tablas de una base de datos;
semiestructurados, como archivos JSON o XML; y no estructurados, como
imágenes, videos, correos electrónicos y publicaciones de redes sociales.

Su objetivo es descubrir patrones, relaciones, tendencias y anomalías que
permitan generar conocimiento útil. Para conseguirlo emplea técnicas como
clasificación, agrupamiento, regresión, correlación, minería de datos y
aprendizaje automático.

\subsection{¿Cuáles son los principales tipos de Big Data Analytics?}

Los principales tipos son:

\begin{enumerate}
    \item \textbf{Analítica descriptiva:} explica qué ocurrió.
    \item \textbf{Analítica diagnóstica:} investiga por qué ocurrió.
    \item \textbf{Analítica predictiva:} estima qué podría ocurrir.
    \item \textbf{Analítica prescriptiva:} recomienda qué acción realizar.
    \item \textbf{Analítica cognitiva:} utiliza inteligencia artificial para
    aprender, interpretar información y relacionarse con las personas.
\end{enumerate}

Estos tipos pueden utilizarse de manera consecutiva. Primero se conoce la
situación actual, después se buscan sus causas, se pronostican resultados y,
finalmente, se recomienda una decisión.

\subsection{¿Qué es la analítica descriptiva?}

La analítica descriptiva estudia datos actuales e históricos para resumir lo
que ocurrió dentro de una organización. Responde principalmente la pregunta:
\textit{¿qué pasó?}

Utiliza operaciones como sumas, promedios, porcentajes, conteos y
comparaciones. Sus resultados normalmente se presentan mediante tablas,
gráficas, indicadores y tableros.

Su función principal es ofrecer una visión clara de la situación. Por ejemplo,
puede mostrar las ventas mensuales, el número de clientes, los productos más
vendidos o el comportamiento de los costos. No explica las causas ni realiza
predicciones.

\subsection{¿Qué es la analítica diagnóstica?}

La analítica diagnóstica analiza los datos para encontrar las razones de un
resultado o problema. Responde la pregunta: \textit{¿por qué ocurrió?}

Utiliza técnicas como correlación, comparación, segmentación, búsqueda
detallada de información y análisis de causa raíz. Por ejemplo, después de
detectar una disminución en las ventas, puede investigar si estuvo relacionada
con los precios, la temporada, la competencia o la falta de inventario.

Se diferencia de la analítica descriptiva porque la descriptiva muestra el
resultado, mientras que la diagnóstica busca sus posibles causas. Debe
recordarse que una correlación no siempre demuestra una relación de causa y
efecto.

\subsection{¿Qué papel tiene la analítica prescriptiva en la toma de decisiones?}

La analítica prescriptiva ayuda a seleccionar la acción más conveniente ante
una situación. Responde preguntas como: \textit{¿qué debemos hacer?} y
\textit{¿cómo podemos obtener el resultado deseado?}

Para generar recomendaciones puede utilizar predicciones, reglas, simulaciones,
optimización matemática y aprendizaje automático. También compara distintos
escenarios y calcula las posibles consecuencias de cada decisión.

Por ejemplo, una empresa puede utilizarla para elegir una ruta de distribución,
definir el nivel adecuado de inventario o determinar la mejor oferta para un
cliente. La decisión final continúa siendo responsabilidad de las personas,
quienes deben considerar el contexto, los riesgos y las limitaciones del
modelo.

\subsection{¿En qué consiste la analítica cognitiva?}

La analítica cognitiva combina Big Data con inteligencia artificial,
aprendizaje automático, aprendizaje profundo, procesamiento del lenguaje
natural y análisis semántico.

Su propósito es procesar información de una manera semejante a ciertas
capacidades humanas. Puede interpretar textos, reconocer imágenes, comprender
preguntas, identificar patrones y aprender de nuevos datos.

Una aplicación cognitiva mejora conforme recibe más información y
retroalimentación. Se utiliza en asistentes virtuales, sistemas de
recomendación, análisis de documentos, atención al cliente y apoyo al
diagnóstico médico. Sus resultados deben ser supervisados porque pueden
contener errores o reproducir sesgos presentes en los datos.

\subsection{Casos de uso de los tipos de Big Data Analytics}

\subsubsection*{Analítica descriptiva: Dow}

La lectura presenta el caso de la empresa química Dow. La organización analizó
información histórica sobre el uso de sus oficinas y laboratorios.

El análisis permitió identificar espacios que se encontraban poco utilizados.
Con esta información, la empresa pudo aprovechar mejor sus instalaciones y
obtener un ahorro aproximado de cuatro millones de dólares por año. Es un caso
descriptivo porque parte de datos históricos para mostrar cómo estaban siendo
utilizados los espacios.

\subsubsection*{Analítica diagnóstica: HelloFresh}

HelloFresh reúne información como ubicación, características demográficas,
preferencias de alimentos y frecuencia de compra.

En el ejemplo de la lectura se detecta un aumento en los pedidos de recetas con
pescado. El análisis diagnóstico busca las características relacionadas con
ese aumento y encuentra una asociación con mujeres que viven en el noreste de
Estados Unidos.

Es importante señalar que este es un caso ilustrativo de correlación. Para
confirmar la causa del comportamiento sería necesario realizar una
investigación adicional.

\subsubsection*{Analítica predictiva: Universidad de Purdue}

La Universidad de Purdue utilizó información académica y de comportamiento
para calcular el nivel de riesgo de sus estudiantes.

El sistema generaba un perfil utilizando colores: verde para una alta
probabilidad de éxito, amarillo para posibles dificultades y rojo para riesgo
de fracaso. Con esta predicción, profesores y estudiantes podían actuar antes
de que el problema aumentara.

Se considera analítica predictiva porque utiliza datos anteriores y actuales
para estimar un posible resultado futuro.

\subsubsection*{Analítica prescriptiva: Aurora Health Care}

En el sector salud, la analítica prescriptiva puede comparar tratamientos,
medicamentos y resultados de pacientes anteriores para recomendar opciones a
los médicos.

La lectura presenta el caso de Aurora Health Care, donde el análisis se utilizó
para apoyar la selección de tratamientos. Además de mejorar las decisiones
médicas, se reportó un ahorro aproximado de seis millones de dólares por año.

La herramienta ofrece recomendaciones, pero el diagnóstico y el tratamiento
final deben ser determinados por profesionales de la salud.

\subsubsection*{Analítica cognitiva: asistentes y plataformas de inteligencia artificial}

Como ejemplos históricos, la lectura menciona Cortana de Microsoft, Siri de
Apple y Watson de IBM.

Estas plataformas pueden procesar lenguaje natural, interpretar solicitudes y
aprender de las interacciones. También pueden analizar documentos, mensajes y
otros datos no estructurados.

Son ejemplos de analítica cognitiva porque combinan información, inteligencia
artificial e interacción con las personas. Algunas de estas tecnologías han
cambiado o han sido sustituidas por asistentes más recientes, pero siguen
siendo ejemplos importantes de la evolución de este tipo de analítica.

\subsection{Principales ventajas de aplicar Big Data Analytics}

\subsubsection*{Administración de riesgos}

El análisis de grandes cantidades de datos ayuda a detectar comportamientos
inusuales, amenazas de seguridad, operaciones fraudulentas y cambios en el
mercado. También permite anticipar algunos riesgos antes de que produzcan
daños importantes.

No elimina completamente la incertidumbre, pero proporciona información para
tomar medidas preventivas.

\subsubsection*{Reducción de costos}

Big Data Analytics permite encontrar procesos lentos, recursos poco utilizados,
errores frecuentes y gastos innecesarios.

Una empresa puede usar estos resultados para optimizar inventarios, reducir
desperdicios, mejorar el mantenimiento de equipos y utilizar de manera más
eficiente su infraestructura. El beneficio depende de que los resultados se
conviertan en acciones reales.

\subsubsection*{Mejores decisiones basadas en datos}

El análisis permite fundamentar las decisiones con información y no solamente
con opiniones. Los datos históricos ayudan a comprender una situación,
mientras que los modelos predictivos permiten prepararse para diferentes
escenarios.

Para que una decisión sea confiable, los datos deben ser correctos, completos,
actuales y relacionados con el problema que se desea resolver.

\subsubsection*{Mejora del desarrollo de productos}

Las organizaciones pueden analizar ventas, opiniones, devoluciones y
preferencias de los clientes para conocer qué características funcionan y
cuáles deben mejorarse.

Esta información ayuda a diseñar productos más cercanos a las necesidades de
los consumidores, personalizar servicios y responder con mayor rapidez a los
cambios de la demanda.

\subsection{Principales retos del Big Data Analytics}

\subsubsection*{Falta de profesionales capacitados}

Los proyectos de Big Data requieren conocimientos de programación,
estadística, bases de datos, infraestructura, seguridad y aprendizaje
automático.

No siempre es fácil encontrar personas que dominen todas estas áreas. Además,
las tecnologías cambian rápidamente, por lo que es necesario mantener una
capacitación continua.

\subsubsection*{Interpretación incorrecta del Big Data}

Un error frecuente es pensar que una mayor cantidad de información siempre
produce mejores resultados. Un conjunto grande, pero incorrecto o poco
relacionado con el problema, puede generar conclusiones equivocadas.

La tecnología tampoco resuelve los problemas por sí sola. Es necesario definir
objetivos claros, seleccionar datos útiles e interpretar correctamente los
resultados. El análisis debe ser un proceso continuo y no una actividad que se
realiza una sola vez.

\subsubsection*{Crecimiento de los datos}

El aumento constante de la información puede superar la capacidad disponible
de almacenamiento, procesamiento y red. También hace más difícil encontrar los
datos que realmente tienen valor.

Conforme crece el volumen, pueden aparecer duplicados, errores,
inconsistencias y costos adicionales. Para enfrentar este reto se necesitan
sistemas escalables, reglas para conservar o eliminar información y controles
de calidad.

\subsubsection*{Selección de herramientas}

Existen muchas herramientas de almacenamiento, procesamiento, visualización e
inteligencia artificial. Cada una tiene características, costos y niveles de
dificultad diferentes.

Una selección incorrecta puede producir gastos innecesarios o una solución
incapaz de crecer. Antes de decidir, la organización debe analizar el volumen
de datos, la velocidad requerida, la experiencia del equipo, la seguridad, la
compatibilidad y el presupuesto. También es recomendable realizar una prueba
pequeña.

\subsubsection*{Seguridad y privacidad}

Concentrar grandes cantidades de información aumenta el impacto de un acceso
no autorizado o una filtración. Además, los datos pueden revelar información
personal incluso cuando esa no era la intención original.

Las organizaciones deben aplicar permisos, cifrado, autenticación, copias de
seguridad, auditorías y políticas de conservación. También deben informar
cómo utilizan los datos, obtener el consentimiento necesario y cumplir las
normas de privacidad aplicables.

La seguridad debe estar presente durante todo el ciclo de vida y no solamente
en la etapa de almacenamiento.

\subsection{Fases del ciclo de vida de Big Data Analytics}

\subsubsection*{Adquisición de datos}

La primera fase consiste en recopilar datos estructurados, semiestructurados y
no estructurados desde distintas fuentes.

Estas fuentes pueden ser bases de datos, sensores, dispositivos móviles,
redes sociales, páginas web, hospitales o sistemas empresariales. Herramientas
como Apache Kafka, Apache Flume, Apache Sqoop y RabbitMQ pueden ayudar a
transportar la información hacia las siguientes etapas.

Antes de recopilar datos es necesario establecer su propósito y verificar que
su obtención sea legal y necesaria.

\subsubsection*{Preprocesamiento}

En esta fase los datos se preparan para que puedan analizarse. La información
recibida puede contener campos vacíos, duplicados, formatos diferentes, errores
o valores fuera de lo normal.

Entre las actividades de preprocesamiento se encuentran:

\begin{itemize}
    \item Integrar información de diferentes fuentes.
    \item Corregir o completar valores faltantes.
    \item Eliminar duplicados.
    \item Convertir variables categóricas en valores adecuados para el modelo.
    \item Escalar o normalizar datos numéricos.
    \item Detectar valores atípicos.
    \item Seleccionar o extraer las características más útiles.
\end{itemize}

Un preprocesamiento correcto mejora la calidad del análisis y reduce la
posibilidad de obtener conclusiones equivocadas.

\subsubsection*{Almacenamiento}

La tercera fase consiste en guardar los datos de forma segura y accesible. El
sistema debe permitir aumentar su capacidad conforme crece la información y
responder con una velocidad adecuada.

Dependiendo del proyecto, pueden utilizarse sistemas distribuidos como HDFS,
bases de datos NoSQL, lagos de datos, almacenes de datos o servicios de nube.

En esta fase también se definen permisos, copias de seguridad, periodos de
conservación y medidas de cifrado.

\subsubsection*{Análisis}

En la última fase se aplican técnicas para descubrir patrones, tendencias,
relaciones y anomalías.

Se pueden utilizar estadísticas, minería de datos, aprendizaje automático,
procesamiento del lenguaje natural y modelos predictivos. Los resultados
deben interpretarse y presentarse mediante reportes, gráficas o tableros
comprensibles.

Después del análisis se toman decisiones y se evalúan sus resultados. Si el
problema no fue resuelto, el ciclo puede repetirse con nuevos datos o con un
método diferente.

\section{Conclusión}

Big Data Analytics transforma grandes cantidades de información en
conocimiento útil. Los diferentes tipos de analítica permiten describir una
situación, investigar sus causas, anticipar resultados, recomendar acciones y
procesar información mediante inteligencia artificial.

Sus beneficios incluyen mejores decisiones, reducción de costos, desarrollo
de productos y administración de riesgos. Sin embargo, para obtener resultados
confiables es necesario contar con datos de calidad, personal capacitado,
herramientas apropiadas y controles de seguridad y privacidad durante todo el
ciclo de vida.